\documentclass[a4paper]{article}

%% Language and font encodings
\usepackage[english]{babel}
\usepackage[utf8]{inputenc}
\usepackage[T1]{fontenc}
\usepackage{amsfonts}

%% Sets page size and margins
\usepackage[a4paper,top=3cm,bottom=2cm,left=3cm,right=3cm,marginparwidth=1.75cm]{geometry}

%% Useful packages
\usepackage{amsmath}
\usepackage{amssymb}
\usepackage{amsthm}
\usepackage{graphicx}
\usepackage[colorinlistoftodos]{todonotes}
\usepackage[colorlinks=true, allcolors=blue]{hyperref}

%% Table packages
\usepackage{array}
\newcolumntype{P}[1]{>{\centering\arraybackslash}p{#1}}
\usepackage{multirow}

\title{Mathematical and Physical Modeling for MATLAB Simulation of Entangled Photon Systems in Optical Fibers}
\author{Marco Cioci}
\date{}

\begin{document}
\maketitle

\tableofcontents

%========================================================
\section{General Framework}

\subsection{Purpose and Scope}

The purpose of this document is to establish a structured working framework for the mathematical modeling and simulation of polarization-entangled photon pairs propagating in optical fibers.

This text is not intended to be part of the final thesis manuscript. Instead, it serves as an internal reference throughout the development process, supporting communication, ensuring consistency of notation, and enabling progressive validation of both theoretical and computational steps.

The primary objective is to organize the physical assumptions, mathematical tools, and simulation strategy in a coherent and reproducible manner, maintaining a direct correspondence between theoretical formulation and MATLAB implementation.

\paragraph{Simulation notes.}
At MATLAB simulator level, this document is treated as the reference specification for notation, model hierarchy, and implementation priorities. The simulator should therefore be organized so that each theoretical block corresponds to a clearly identifiable module, making it possible to validate intermediate results step by step and to extend the code without changing the whole structure.

\subsection{Theoretical Formalism and Notation}

The system is described within the formalism of quantum information theory. The relevant mathematical structures include Hilbert spaces, quantum states, tensor products, unitary operators, density operators, and quantum measurements. Classical optics is introduced only when necessary to motivate the physical origin of the effective quantum transformations.

The notation and formal tools are aligned, as far as possible, with the framework of \textit{Quantum Computation and Quantum Information} by Nielsen and Chuang. In particular, the computational basis of a single qubit is identified with the polarization basis:
\[
|0\rangle \equiv |H\rangle,
\qquad
|1\rangle \equiv |V\rangle,
\]
where \( |H\rangle \) and \( |V\rangle \) denote horizontal and vertical polarization states.

This choice fixes from the beginning the language in which the system is described: polarization states are treated as qubit states, bipartite photon pairs are treated as two-qubit systems, and optical transformations are represented by operators acting on finite-dimensional Hilbert spaces.

\paragraph{Simulation notes.}
The simulator should adopt from the beginning a quantum-information-oriented representation: state vectors for pure states, density matrices for mixed states, tensor-product composition for bipartite systems, and matrix operators for evolutions and measurements. The polarization encoding must be fixed once and for all as \( |H\rangle \leftrightarrow |0\rangle \), \( |V\rangle \leftrightarrow |1\rangle \), so that all MATLAB routines use the same basis convention.

At the current implementation stage, this convention is already enforced consistently in the codebase. In particular, Bell states, local fiber unitaries, bipartite tensor-product operators, and correlation observables are all represented in the ordered computational basis
\[
\{|00\rangle,|01\rangle,|10\rangle,|11\rangle\}.
\]

\subsection{Physical Scenario and Experimental Setting}

The physical setting is as follows. A spontaneous parametric down-conversion (SPDC) source generates pairs of polarization-entangled photons, ideally prepared in a Bell state. The two photons are injected into distinct optical paths, denoted \(A\) and \(B\), and propagate through independent optical fibers.

Each fiber acts on the polarization degree of freedom of the corresponding photon. In the most general description, this action is represented by a unitary operator on a single-qubit Hilbert space. However, as a first approximation motivated by experimental considerations, the dominant residual effect is modeled as a phase rotation around the computational \(Z\)-axis of the Bloch sphere.

The aim is therefore not to describe the full electromagnetic propagation problem from first principles at the initial stage, but rather to identify an effective quantum model that captures the physically relevant action of the fibers on polarization entanglement.

\paragraph{Simulation notes.}
At simulator level, the physical picture is reduced to the minimum set of objects needed for computation: an initial entangled two-qubit state, two labeled channels \(A\) and \(B\), and one local transformation per channel. The first implementation should neglect unnecessary optical detail and describe each fiber only through its effective action on polarization.

In the presently implemented experiment, the local evolution is chosen as
\[
U_A = I,
\qquad
U_B = U(\theta),
\]
so that the full state evolution is generated by a single swept phase parameter. This provides the first operational validation setting for the simulator.

%========================================================
\section{Modeling Strategy}

\subsection{Progressive Modeling Approach}

The modeling approach follows a progressive refinement strategy, in which the system is described through successive levels of approximation of increasing physical complexity.

Each stage introduces additional effects while remaining consistent with the previous ones. This ensures both analytical transparency and modularity in the implementation.

From a theoretical perspective, this approach isolates the contribution of individual physical mechanisms. From a computational perspective, it enables a layered structure in the simulator, where each refinement corresponds to an additional module.

\paragraph{Simulation notes.}
The MATLAB code should be designed incrementally. A working minimal simulator must exist before more realistic effects are added. Each new level should extend the previous one rather than replace it, so that one can compare simplified and refined models under controlled conditions and verify where each physical effect begins to matter.

\subsection{Hierarchy of Modeling Levels}

The overall workflow can be summarized as:
\[
\begin{array}{c}
\text{Ideal Bell State} \\
\downarrow \\
\text{Local Phase Model (Z-rotations)} \\
\downarrow \\
\text{General Unitary Transformations} \\
\downarrow \\
\text{Mixed States (Density Operators)} \\
\downarrow \\
\text{Noisy Quantum Channels} \\
\downarrow \\
\text{Full Fiber Models}
\end{array}
\]

More explicitly:

\begin{itemize}
    \item \textbf{Level 1: Minimal model} \\
    Pure, maximally entangled state. Fiber action reduced to a relative phase. Analytical dependence of correlations on a single parameter.

    \item \textbf{Level 2: General unitary dynamics} \\
    Arbitrary polarization rotations described by local unitary operators.

    \item \textbf{Level 3: Mixed-state description} \\
    Introduction of density operators to model decoherence and loss of purity.

    \item \textbf{Level 4: Noisy channels} \\
    Effective quantum channels (e.g.\ depolarizing maps) to represent fiber-induced noise.

    \item \textbf{Level 5: Physical fiber models} \\
    More detailed descriptions, such as concatenations of random unitary segments or stochastic models of polarization evolution.
\end{itemize}

This hierarchy provides both a conceptual roadmap and an implementation roadmap. The first levels are designed to remain analytically tractable and physically interpretable, while the later levels are intended to increase realism and to connect more directly with experimentally relevant imperfections.

\paragraph{Simulation notes.}
This hierarchy should be mirrored directly in the codebase. A convenient choice is to implement one MATLAB function or class layer for each level: state preparation, local unitary evolution, density-matrix conversion, channel application, and advanced fiber propagation. This makes later thesis work easier, because each refinement corresponds to a visible computational step.

At present, the implementation covers the first level completely and has already prepared part of the structural basis for later extensions. In particular, the current code already distinguishes state preparation routines, local propagation models, correlation-evaluation routines, validation utilities, and experiment scripts.

\subsection{Baseline Model: Pure-State Description}

In the first stage, the global bipartite system is described as a pure state in \(\mathbb{C}^2 \otimes \mathbb{C}^2\). Fiber propagation is modeled through local unitary transformations, with particular focus on relative phase shifts.

The main objective is to analyze how these phase shifts affect measurable two-photon correlations. The relevant observables are operators of the form
\[
\sigma_i \otimes \sigma_j,
\qquad i,j \in \{x,y,z\},
\]
as well as observables corresponding to arbitrary measurement directions on the Bloch sphere.

This baseline model is the natural entry point because it is simple enough to allow explicit formulas, direct visualization of the state dependence on propagation parameters, and immediate comparison between analytical and numerical results.

\paragraph{Simulation notes.}
The first operational simulator should work entirely with pure two-qubit states represented as \(4 \times 1\) complex vectors. The initial tasks are: generate the input Bell state, apply local matrices \(U_A\) and \(U_B\), and compute expectation values of selected observables. This level should already produce plots or tables showing how correlations vary with the relative phase.

This stage is now concretely represented in MATLAB by a set of routines for Bell-state generation, tensor-product local evolution, Pauli-operator construction, expectation-value evaluation, and correlation extraction. In the current version, the standard aligned correlations
\[
\langle \sigma_x \otimes \sigma_x \rangle,\qquad
\langle \sigma_y \otimes \sigma_y \rangle,\qquad
\langle \sigma_z \otimes \sigma_z \rangle
\]
are explicitly computed and can be plotted as functions of the swept phase parameter.

\subsection{Extension to Mixed-State Formalism}

At a more advanced stage, the model incorporates decoherence and depolarization effects. In this regime, the pure-state description becomes insufficient, and density operators are required.

This framework allows one to describe:
\begin{itemize}
    \item reduced states of individual photons,
    \item loss of coherence due to frequency-dependent propagation,
    \item noisy quantum channels acting on polarization.
\end{itemize}

It also enables the computation of quantities such as fidelity, concurrence, and other entanglement measures.

The introduction of mixed states is therefore not a merely formal refinement. It is required whenever one wishes to describe partial information, environmental coupling, statistical mixtures, or effective non-unitary dynamics arising from unresolved physical degrees of freedom.

\paragraph{Simulation notes.}
Once the pure-state stage is stable, the simulator should be extended so that every relevant routine can also accept density matrices. In practice, this means introducing \(4 \times 4\) matrix representations, partial trace routines, noise-channel application, and metric-evaluation functions. The pure-state code should remain available, but it should become a special case of the density-operator framework whenever possible.

\subsection{Simulator Design and Architecture}

The simulator is designed as a modular implementation of the theoretical model.

The initial modules will:
\begin{itemize}
    \item generate ideal Bell states,
    \item apply local fiber transformations,
    \item compute correlation functions.
\end{itemize}

Subsequent modules will:
\begin{itemize}
    \item introduce noise channels,
    \item compute reduced density operators,
    \item evaluate observables and entanglement measures,
    \item extract experimentally relevant quantities such as Bell parameters and Stokes-like expectation values.
\end{itemize}

The theoretical development is therefore directly coupled to the computational structure, ensuring consistency, flexibility, and traceability of results.

The architecture should distinguish clearly between reusable physical kernels and higher-level experiment scripts. In this way the simulator can evolve from a set of exploratory routines into a compact and reusable research tool.

\paragraph{Simulation notes.}
The architecture should be kept explicitly modular: one block for state preparation, one for channel action, one for observables, one for diagnostics, and one for post-processing. At MATLAB level, this suggests separating reusable numerical kernels from scripts used only for running specific experiments, so that the simulator can later become a small reusable library.

At the current stage, this modular principle is already reflected in the directory organization:
\begin{itemize}
    \item \texttt{Core}: state vectors, operators, expectation values, and correlation routines;
    \item \texttt{Models}: physical propagation models such as the reduced phase unitary;
    \item \texttt{Utils}: validation routines, tensor-product wrapper, and plotting helpers;
    \item \texttt{Experiments}: scripts or functions implementing complete studies, such as phase sweeps.
\end{itemize}
This separation is not only cosmetic: it enforces the distinction between general mathematical building blocks and physically specific model choices.

%========================================================
\section{Physical and Mathematical Representation}

\subsection{Quantum System and Hilbert Space Structure}

The starting point of the model is a source of polarization-entangled photon pairs generated by spontaneous parametric down-conversion (SPDC), in a configuration in which the two output paths are denoted \(A\) and \(B\). When the relevant optical alternatives are made indistinguishable, the source can ideally prepare one of the Bell states of two polarization qubits.

In the present work, the natural encoding of the qubit is the polarization degree of freedom:
\[
|H\rangle \equiv |0\rangle,
\qquad
|V\rangle \equiv |1\rangle.
\]
Accordingly, the computational basis of the bipartite system is
\[
\{
|00\rangle,\,
|01\rangle,\,
|10\rangle,\,
|11\rangle
\},
\]
that is,
\[
\{
|H\rangle_A |H\rangle_B,\,
|H\rangle_A |V\rangle_B,\,
|V\rangle_A |H\rangle_B,\,
|V\rangle_A |V\rangle_B
\}.
\]

The Hilbert space of the total system is therefore
\[
\mathcal{H}_{AB} = \mathcal{H}_A \otimes \mathcal{H}_B
\cong \mathbb{C}^2 \otimes \mathbb{C}^2,
\]
which has dimension \(4\). As a consequence, a pure state of the two-photon system is represented by a complex vector with four components in the computational basis, while a general linear operator acting on the full system is represented by a \(4 \times 4\) matrix.

This representation fixes the elementary mathematical objects used throughout the model: single-photon polarization states are qubit vectors in \(\mathbb{C}^2\), two-photon states live in the tensor-product space \(\mathbb{C}^2 \otimes \mathbb{C}^2\), and observables or evolutions are represented by matrices acting on these spaces.

\paragraph{Simulation notes.}
At MATLAB level, the natural internal representation is therefore immediate: a bipartite pure state is a complex vector of length \(4\), and bipartite operators are \(4 \times 4\) matrices. The simulator should store basis ordering explicitly as
\[
\{|00\rangle,|01\rangle,|10\rangle,|11\rangle\},
\]
since all subsequent tensor products, observables, and density matrices depend on this convention.

This basis ordering is now a concrete implementation convention of the codebase. Tensor-product operators, Bell-state vectors, local unitary evolutions, and plotted correlation outputs all rely on this ordering and should therefore continue to use it consistently in every future extension.

\subsection{Bell-State Initialization}

As a reference initial condition, we consider the Bell state
\[
|\Psi^{+}\rangle
=
\frac{1}{\sqrt{2}}
\left(
|H\rangle_A |V\rangle_B
+
|V\rangle_A |H\rangle_B
\right),
\]
which, in computational notation, is written as
\[
|\Psi^{+}\rangle
=
\frac{1}{\sqrt{2}}
\left(
|01\rangle + |10\rangle
\right).
\]

In the ordered computational basis
\[
\{|00\rangle, |01\rangle, |10\rangle, |11\rangle\},
\]
this state is represented by the column vector
\[
|\Psi^{+}\rangle
=
\frac{1}{\sqrt{2}}
\begin{pmatrix}
0\\
1\\
1\\
0
\end{pmatrix}.
\]

This state is maximally entangled and constitutes the natural input of the first version of the simulator. It is also the reference state with respect to which quantities such as fidelity and Bell-inequality violation will later be evaluated.

The choice of a Bell state is natural because it provides both a clean theoretical benchmark and a physically relevant entangled resource. It also allows the effect of the fibers to be interpreted directly in terms of changes in entanglement, coherence, and two-photon correlations.

\paragraph{Simulation notes.}
The simulator should include a dedicated state-preparation routine returning the Bell vector in canonical basis ordering. This reference state should be used systematically for debugging, verification of normalization, comparison with analytical formulas, and later evaluation of fidelity after propagation through the fibers.

At the current stage, the simulator already includes a Bell-state generator able to return the four standard Bell states
\[
|\Psi^{\pm}\rangle
=
\frac{1}{\sqrt{2}}\left(|01\rangle \pm |10\rangle\right),
\qquad
|\Phi^{\pm}\rangle
=
\frac{1}{\sqrt{2}}\left(|00\rangle \pm |11\rangle\right),
\]
although the present baseline experiment uses \( |\Psi^+\rangle \) as the default reference input.

\subsection{Fiber Propagation as Local Unitary Evolution}

Each photon propagates through its own optical fiber. At the level of the polarization qubit, the action of each fiber is described by a linear operator acting locally on the corresponding single-photon Hilbert space. In the most general ideal lossless description, such an operator is unitary. Therefore, if \(U_A\) and \(U_B\) denote the single-qubit operators associated with the two fibers, the evolution of the initial two-photon state is
\[
|\psi_{\mathrm{out}}\rangle
=
(U_A \otimes U_B)\, |\psi_{\mathrm{in}}\rangle.
\]

This expression is the natural bipartite extension of single-qubit evolution and follows directly from the tensor-product structure of composite quantum systems. It provides the correct starting point for any subsequent refinement of the model, including random unitary perturbations, concatenated fiber sections, or effective channel descriptions.

In this formulation, the two fibers act independently on the two polarization subsystems, and the total action is obtained by tensor composition of the corresponding local transformations. This is the mathematically appropriate description whenever the fibers do not generate direct interaction between the two photons.

\paragraph{Simulation notes.}
The fundamental evolution primitive in MATLAB should be the tensor-product action
\[
\psi_{\mathrm{out}} = (U_A \otimes U_B)\psi_{\mathrm{in}}.
\]
It is convenient to isolate this in a single function, so that later changes in the single-fiber model automatically propagate to the full bipartite simulation.

This abstraction is already reflected in the present implementation, where a dedicated routine applies local single-qubit matrices to a bipartite pure state by constructing the global operator \(U_A \otimes U_B\). The tensor-product operation itself is also isolated in a dedicated wrapper function, so that the notation used in code remains aligned with the mathematical notation used in the model.

\subsection{Reduced Phase Model}

Although the most general fiber action may involve arbitrary polarization rotations, birefringence, and mode-coupling effects, the first approximation adopted here is simpler and is directly motivated by the experimental control scheme. After partial compensation of polarization drifts, the dominant residual effect is modeled as a relative phase shift between the horizontal and vertical components, that is, as an effective rotation around the \(Z\)-axis of the Bloch sphere.

At the single-qubit level, this is represented by a diagonal unitary of the form
\[
U(\theta)
=
\begin{pmatrix}
1 & 0\\
0 & e^{i\theta}
\end{pmatrix},
\]
up to an irrelevant global phase. If the two fibers induce phases \(\theta_A\) and \(\theta_B\), then the total transformation is
\[
U_A \otimes U_B
=
U(\theta_A)\otimes U(\theta_B).
\]

Applying this operator to the reference Bell state \( |\Psi^+\rangle \) gives
\[
|\psi(\theta_A,\theta_B)\rangle
=
\frac{1}{\sqrt{2}}
\left(
e^{i\theta_B}|01\rangle
+
e^{i\theta_A}|10\rangle
\right).
\]
By factoring out a global phase, which has no physical effect on measurement statistics, this state can be rewritten as
\[
|\psi(\theta)\rangle
=
\frac{1}{\sqrt{2}}
\left(
|01\rangle
+
e^{i\theta}|10\rangle
\right),
\]
where
\[
\theta = \theta_A - \theta_B.
\]

Thus, in the first effective model, the action of the two fibers is reduced to a single physically relevant parameter: the relative phase \(\theta\) between the two components of the entangled state. This parameter is expected to depend on propagation conditions such as birefringence, path length, and environmental fluctuations, and it will be the main control variable in the first stage of the simulator.

This reduced description is especially useful because it isolates the effect of interest while preserving the essential entangled two-qubit structure. It is therefore a controlled approximation rather than an arbitrary simplification.

\paragraph{Simulation notes.}
This is the first concrete propagation model to implement. The simulator should expose \(\theta_A\) and \(\theta_B\) as input parameters, but internally recognize that the observable behavior depends only on \(\theta=\theta_A-\theta_B\). In practice, it is useful to keep both forms available: the two-phase description for physical interpretation, and the reduced single-parameter description for plots, parameter sweeps, and analytical checks.

At the current implementation stage, the reduced model is realized through a dedicated single-qubit phase-unitary routine and is used operationally in the first complete experiment by fixing one arm to the identity and sweeping the phase in the other arm. This realizes a concrete numerical study of the one-parameter family \( |\psi(\theta)\rangle \).

\subsection{Interpretation and Role of the Reduced Model}

The reduced phase model is particularly convenient for three reasons. First, it preserves the pure-state description, allowing the initial implementation to remain simple and analytically transparent. Second, it isolates the effect that is most directly relevant for the first experimental analysis, namely the deformation of two-photon correlations induced by a relative phase. Third, it prepares the transition to more general models, since the same bipartite formalism can later be extended to density operators and noisy channels without changing the conceptual structure of the problem.

For this reason, the first stage of the simulator will be built around the family of states
\[
|\psi(\theta)\rangle
=
\frac{1}{\sqrt{2}}
\left(
|01\rangle + e^{i\theta}|10\rangle
\right),
\]
from which one will compute expectation values of local and nonlocal observables, correlation functions, and Bell-type quantities as functions of \(\theta\).

The role of this model is therefore twofold. On the one hand, it is a physically motivated approximation to the propagation problem. On the other hand, it is the minimal fully operational model from which one can derive and test measurable predictions in a clean and controlled way.

\paragraph{Simulation notes.}
The first validated MATLAB study should therefore be a parameter sweep in \(\theta\), with numerical evaluation of state vectors, expectation values, and correlation quantities. This gives a minimal but complete simulator: one input family of states, one effective propagation parameter, and a set of measurable outputs that can be directly plotted and compared against theory.

This step is now represented concretely by a sweep routine returning a structured collection of results containing the sampled phase values and the corresponding correlation outputs. A plotting utility then allows one to visualize selected correlations directly as functions of \(\theta\), making the reduced model not only conceptually defined but also already operational at code level.

\subsection{Correlation Observables: Theory and Implementation}

At the present stage of development, the observables explicitly implemented in the simulator are the standard aligned two-qubit Pauli correlations:
\[
\sigma_x \otimes \sigma_x,\qquad
\sigma_y \otimes \sigma_y,\qquad
\sigma_z \otimes \sigma_z.
\]

These are constructed from the single-qubit Pauli matrices
\[
\sigma_x=
\begin{pmatrix}
0 & 1\\
1 & 0
\end{pmatrix},
\qquad
\sigma_y=
\begin{pmatrix}
0 & -i\\
i & 0
\end{pmatrix},
\qquad
\sigma_z=
\begin{pmatrix}
1 & 0\\
0 & -1
\end{pmatrix}.
\]

Given a pure two-qubit state \( |\psi\rangle \), the expectation values are defined as
\[
\langle \sigma_i \otimes \sigma_i \rangle
=
\langle \psi|
(\sigma_i \otimes \sigma_i)
|\psi\rangle,
\qquad i\in\{x,y,z\}.
\]

\paragraph{Analytical behavior under the reduced phase model.}

Within the reduced phase model, the initial Bell state
\[
|\psi^+\rangle = \frac{|01\rangle + |10\rangle}{\sqrt{2}}
\]
evolves, after propagation in the fiber, into
\[
|\psi(\theta)\rangle
=
\frac{|01\rangle + e^{i\theta}|10\rangle}{\sqrt{2}}.
\]

For this family of states, the correlation observables can be computed analytically.

\begin{itemize}
    \item \textbf{Correlation \( \sigma_x \otimes \sigma_x \):}
    \[
    \langle \sigma_x \otimes \sigma_x \rangle
    =
    \cos\theta.
    \]

    \item \textbf{Correlation \( \sigma_y \otimes \sigma_y \):}
    \[
    \langle \sigma_y \otimes \sigma_y \rangle
    =
    \cos\theta.
    \]

    \item \textbf{Correlation \( \sigma_z \otimes \sigma_z \):}
    \[
    \langle \sigma_z \otimes \sigma_z \rangle
    =
    -1.
    \]
\end{itemize}

These results show that:
\begin{itemize}
    \item The fiber-induced phase \(\theta\) affects only correlations in the transverse plane (\(x,y\)).
    \item The longitudinal correlation (\(z\)) remains invariant, reflecting the fact that the relative phase does not alter populations in the computational basis.
\end{itemize}

This structure is a direct consequence of the fact that the fiber, in the reduced model, acts as a local unitary equivalent to a rotation around the \(z\)-axis in the Bloch sphere representation.

\paragraph{Simulation notes.}

At implementation level, the simulator contains routines for:
\begin{itemize}
    \item generating the Pauli matrices,
    \item constructing tensor-product operators,
    \item evaluating expectation values via inner products,
    \item packaging correlation outputs into structured data containers.
\end{itemize}

Additional helper functions verify:
\begin{itemize}
    \item normalization of input state vectors,
    \item unitarity of operators used to model fiber evolution.
\end{itemize}

These observables constitute the primary diagnostic outputs of the simulator and provide a direct bridge between theoretical predictions and numerical results. Their dependence on \(\theta\) enables a clear validation pipeline through comparison with simulated data.

\begin{figure}[h]
    \centering
    \includegraphics[width=0.7\linewidth]{../Images/correlation_observables_theta.png}
    \caption{
Numerical evaluation of the two-qubit correlation observables 
$c_{xx}$, $c_{yy}$, and $c_{zz}$ as functions of the relative phase $\theta$ 
introduced by the optical fiber.
The results are obtained by applying the fiber-induced unitary evolution to the initial Bell state 
and computing the expectation values of the corresponding Pauli operators.}
    \label{fig:correlations_numerical}
\end{figure}

\begin{figure}[h]
    \centering
    \includegraphics[width=0.7\linewidth]{../Images/experiment_1_theory_vs_experiment.png}
    \caption{
Comparison between numerical simulation and analytical predictions for the two-qubit correlation observables 
$c_{xx}$, $c_{yy}$, and $c_{zz}$ as functions of $\theta$.
Starting from the evolved state 
$\frac{1}{\sqrt{2}} (|01\rangle + e^{i\theta} |10\rangle)$, 
the analytical model predicts 
$c_{xx} = c_{yy} = \cos(\theta)$ and $c_{zz} = -1$. 
The numerical results are in full agreement with these expressions. 
Solid lines denote theoretical predictions, while markers correspond to simulated values.}
    \label{fig:correlations_validation}
\end{figure}

\section{Numerical Experiments and Results}

\subsection{Experiment 1: Deterministic Phase Sweep}

\subsubsection{Objective}

The first experiment investigates the effect of a controlled relative phase $\theta$ on a maximally entangled two-qubit state. The goal is to characterize how the correlation observables evolve under a unitary phase shift induced by the fiber.

\subsubsection{Model}

The initial Bell state is taken as
\[
|\psi_+\rangle = \frac{1}{\sqrt{2}} \left( |01\rangle + |10\rangle \right).
\]

After propagation through the fiber, the system evolves into
\[
|\psi(\theta)\rangle = \frac{1}{\sqrt{2}} \left( |01\rangle + e^{i\theta} |10\rangle \right).
\]

The following correlation observables are evaluated:
\[
c_{xx} = \langle \sigma_x \otimes \sigma_x \rangle, \quad
c_{yy} = \langle \sigma_y \otimes \sigma_y \rangle, \quad
c_{zz} = \langle \sigma_z \otimes \sigma_z \rangle.
\]

\subsubsection{Analytical Predictions}

A direct calculation yields:
\[
c_{xx} = \cos(\theta), \quad
c_{yy} = \cos(\theta), \quad
c_{zz} = -1.
\]

This result highlights that the phase $\theta$ only affects correlations in the transverse ($x,y$) planes, while the $z$-correlation remains invariant.

\subsubsection{Numerical Results}

The simulation performs a sweep over $\theta \in [0, 2\pi]$, computing the correlation observables for each value.

\begin{figure}[h]
    \centering
    \includegraphics[width=0.75\linewidth]{../Images/correlation_observables_theta.png}
    \caption{Correlation observables $c_{xx}, c_{yy}, c_{zz}$ as functions of the relative phase $\theta$.}
\end{figure}

\subsubsection{Discussion}

The numerical results perfectly match the analytical predictions. In particular:

\begin{itemize}
    \item $c_{xx}$ and $c_{yy}$ exhibit a cosine dependence on $\theta$
    \item $c_{zz}$ remains constant at $-1$
\end{itemize}

This confirms that the fiber, under the simplified model, acts as a relative phase shifter that preserves maximal entanglement.

\bigskip

\subsection{Experiment 2: Random Phase Ensemble}

\subsubsection{Objective}

The second experiment extends the model by introducing randomness in the phase $\theta$, simulating realistic fiber conditions where birefringence fluctuations induce stochastic phase shifts.

Instead of a single pure state, the system is described by an ensemble of states:
\[
\{ |\psi(\theta_k)\rangle \}_{k=1}^{N}
\]

The goal is to evaluate how statistical averaging affects observable correlations.

\subsubsection{Ensemble Model}

Each realization is given by
\[
|\psi(\theta_k)\rangle = \frac{1}{\sqrt{2}} \left( |01\rangle + e^{i\theta_k} |10\rangle \right),
\]

and the ensemble density matrix is defined as
\[
\rho_{\text{ensemble}} = \frac{1}{N} \sum_{k=1}^{N} |\psi(\theta_k)\rangle \langle \psi(\theta_k)|.
\]

\subsubsection{Analytical Predictions}

The ensemble-averaged correlations are given by:
\[
c_{xx} = \langle \cos(\theta) \rangle, \quad
c_{yy} = \langle \cos(\theta) \rangle, \quad
c_{zz} = -1,
\]

where the average is taken over the phase distribution.

\subsubsection{Numerical Results}

Three different distributions for $\theta$ are considered:

\begin{itemize}
    \item Constant phase
    \item Uniform distribution
    \item Gaussian distribution
\end{itemize}

\begin{figure}[h]
    \centering
    \includegraphics[width=0.9\linewidth]{../Images/Experiment_2/experiment_2_constant_distribution.png}
    \caption{Experiment 2 results for constant phase distribution.}
\end{figure}

\begin{figure}[h]
    \centering
    \includegraphics[width=0.9\linewidth]{../Images/Experiment_2/experiment_2_uniform_distribution.png}
    \caption{Experiment 2 results for uniform phase distribution.}
\end{figure}

\begin{figure}[h]
    \centering
    \includegraphics[width=0.9\linewidth]{../Images/Experiment_2/experiment_2_normal_distribution.png}
    \caption{Experiment 2 results for Gaussian phase distribution.}
\end{figure}

\subsubsection{Discussion}

The results show a clear transition from coherent to partially decohered behavior:

\begin{itemize}
    \item For constant $\theta$, the system behaves identically to Experiment 1
    \item For distributed $\theta$, $c_{xx}$ and $c_{yy}$ decrease in magnitude
    \item For broad distributions, $\langle \cos(\theta) \rangle \to 0$, leading to vanishing transverse correlations
\end{itemize}

This effect corresponds to a loss of phase coherence induced by the fiber.

\subsubsection{Physical Interpretation}

Although each realization remains a pure maximally entangled state, the ensemble average produces a mixed state. This reflects the physical situation where phase fluctuations occur on timescales shorter than the measurement resolution.

Importantly:

\begin{itemize}
    \item Entanglement is not necessarily destroyed at the single-shot level
    \item However, observable correlations are reduced due to phase averaging
\end{itemize}

This behavior represents the first step toward modeling decoherence effects in optical fibers.


%========================================================
\section{References}

\begin{itemize}
    \item M. A. Nielsen and I. L. Chuang, \textit{Quantum Computation and Quantum Information}, Cambridge University Press.
\end{itemize}

\end{document}